\section*{Spivak Capítulo 2, Problema 4}
\addcontentsline{toc}{section}{Spivak Capítulo 2, Problema 4}

\subsection*{(a) Identidad de Vandermonde}

	extbf{Demostrar que:}
\[
\sum_{k=0}^{\ell} \binom{n}{k}\binom{m}{\ell-k} = \binom{n+m}{\ell} \qquad\text{(Identidad de Vandermonde)}
\]

	extbf{Demostración:}

Comparando coeficientes en la identidad
\[
(1+x)^n(1+x)^m=(1+x)^{n+m}
\]
obtenemos la igualdad anunciada. Al expandir ambos factores con el teorema binomial, el coeficiente de $x^{\ell}$ en el producto es
\[
\sum_{k=0}^{\ell}\binom{n}{k}\binom{m}{\ell-k},
\]
mientras que en el otro miembro dicho coeficiente es $\binom{n+m}{\ell}$.

También hay una prueba puramente combinatoria: elegir $\ell$ elementos de un conjunto dividido en dos bloques (de tamaños $n$ y $m$) equivale a fijar cuántos se eligen del primer bloque ($k$) y cuántos del segundo ($\ell-k$), y sumar sobre $k$ da todas las elecciones.

\subsection*{(b) Suma de Cuadrados de Binomios}

\textbf{Demostrar que:}
\[
\sum_{k=0}^{n} \binom{n}{k}^2 = \binom{2n}{n}
\]

\textbf{Ayuda:} Aplicar el teorema binomial a $(1+x)^n(1+x)^n$.

\textbf{Demostración:}

Consideramos la identidad algebraica:
\[
(1+x)^n \cdot (1+x)^n = (1+x)^{2n}
\]

Expandimos el lado izquierdo usando el teorema binomial:
\[
\left(\sum_{k=0}^{n} \binom{n}{k} x^k\right) \left(\sum_{j=0}^{n} \binom{n}{j} x^j\right)
\]

Al multiplicar estas series, el coeficiente de $x^n$ es:
\[
\sum_{k=0}^{n} \binom{n}{k} \binom{n}{n-k}
\]

Pero $\binom{n}{n-k} = \binom{n}{k}$ por simetría, luego:
\[
\sum_{k=0}^{n} \binom{n}{k}^2
\]

Por el lado derecho, la expansión binomial de $(1+x)^{2n}$ es:
\[
(1+x)^{2n} = \sum_{m=0}^{2n} \binom{2n}{m} x^m
\]

El coeficiente de $x^n$ en esta expansión es simplemente $\binom{2n}{n}$.

Como ambos lados de la igualdad $(1+x)^{2n} = (1+x)^n(1+x)^n$ representan el mismo polinomio, los coeficientes de $x^n$ deben ser iguales:
\[
\sum_{k=0}^{n} \binom{n}{k}^2 = \binom{2n}{n}
\]

\qed

\textbf{Interpretación Combinatoria:}

El lado derecho $\binom{2n}{n}$ cuenta el número de formas de elegir $n$ objetos de un conjunto de $2n$ objetos.

El lado izquierdo $\sum_{k=0}^{n} \binom{n}{k}^2$ cuenta lo mismo de una forma diferente: particionamos los $2n$ objetos en dos grupos de $n$ cada uno, y luego elegimos $k$ del primer grupo y $n-k$ del segundo grupo. Sumando sobre todos los posibles valores de $k$, obtenemos todas las formas de elegir $n$ objetos totales.

\subsection*{Ejemplos}

\paragraph{Ejemplo 1 (Vandermonde).} Tomemos $n=5,m=3,\ell=4$. Calculamos
\[
\sum_{k=0}^{4}\binom{5}{k}\binom{3}{4-k} = \binom{5}{1}\binom{3}{3}+\binom{5}{2}\binom{3}{2}+\binom{5}{3}\binom{3}{1}+\binom{5}{4}\binom{3}{0}
\]
\[
=5\cdot1+10\cdot3+10\cdot3+5\cdot1=5+30+30+5=70.
\]
Y efectivamente $\binom{8}{4}=70$.

\paragraph{Ejemplo 2 (Suma de cuadrados).} Tomemos $n=4$. Entonces
\[
\sum_{k=0}^{4}\binom{4}{k}^2=1^2+4^2+6^2+4^2+1^2=1+16+36+16+1=70,
\]
y $\binom{8}{4}=70$, coincidiendo con la identidad del apartado (b).
